\documentclass{ximera}
%verbeterd 28/10/26
\title{Vraagstukken positie, snelheid en versnelling: moeilijke oefeningen}
\author{Daan Lipkens}

\begin{document}
\begin{abstract}
    Moeilijkere oefeningen over positie, snelheid en versnelling bij rechtlijnige en kromlijnige bewegingen, van makkelijk naar moeilijk.
\end{abstract}
\maketitle

\section{Moeilijke oefeningen: positie, snelheid en versnelling}

%———————————————————————————————————————
% OEFENING 1 (makkelijkst): verplaatsing, afgelegde weg en gemiddelde snelheid
\begin{exercise}
    Een koerier vertrekt in het depot ($x = 0{,}0 \, \mathrm{km}$) en rijdt over een rechte weg langs de $x$-as. Eerst rijdt hij $24 \, \mathrm{km}$ in de positieve zin met $48 \, \mathrm{km/h}$. Daarna wacht hij $15 \, \mathrm{min}$ bij een klant. Tot slot rijdt hij $6{,}0 \, \mathrm{km}$ terug in de negatieve zin met $36 \, \mathrm{km/h}$.

    \begin{question}
        Bereken de verplaatsing $\Delta x$ en de afgelegde weg $s$ over de volledige rit.

        \[ \Delta x = \answer[tolerance=0.01,onlineshowanswerbutton]{18} \, \mathrm{km} \qquad \text{en} \qquad s = \answer[tolerance=0.01,onlineshowanswerbutton]{30} \, \mathrm{km} \]

        \begin{solution}\nl
            \underline{Gegevens:}
            \begin{align*}
                & x_{\text{begin}} = 0{,}0 \, \mathrm{km} & & \text{heen: } 24 \, \mathrm{km} \text{ in de positieve zin}\\
                & \text{terug: } 6{,}0 \, \mathrm{km} \text{ in de negatieve zin} & &
            \end{align*}

            \underline{Gevraagd:} $\Delta x$ en $s$\\

            \underline{Oplossing:}\\
            De eindpositie is $x_{\text{eind}} = 24 \, \mathrm{km} - 6{,}0 \, \mathrm{km} = 18 \, \mathrm{km}$. De verplaatsing is het verschil tussen eind- en beginpositie:
            \[
                \Delta x = x_{\text{eind}} - x_{\text{begin}} = 18 \, \mathrm{km} - 0{,}0 \, \mathrm{km} = 18 \, \mathrm{km}
            \]
            De afgelegde weg is de volledige afgelegde afstand, ongeacht de zin:
            \[
                s = 24 \, \mathrm{km} + 6{,}0 \, \mathrm{km} = 30 \, \mathrm{km}
            \]
        \end{solution}
    \end{question}

    \begin{question}
        Bereken de gemiddelde snelheid $\langle v_x \rangle$ over de volledige rit, inclusief de wachttijd, in $\mathrm{m/s}$ en in $\mathrm{km/h}$.

        \[ \langle v_x \rangle = \answer[tolerance=0.01,onlineshowanswerbutton]{5.45} \, \mathrm{m/s} = \answer[tolerance=0.01,onlineshowanswerbutton]{19.6} \, \mathrm{km/h} \]

        \begin{hint}
            Je hebt de totale duur van de rit nodig. Welke drie tijdsintervallen moet je optellen?
        \end{hint}

        \begin{solution}\nl
            \underline{Gegevens:}
            \begin{align*}
                & \Delta x = 18 \, \mathrm{km} = 18 \cdot 10^{3} \, \mathrm{m} & & v_{\text{heen}} = 48 \, \mathrm{km/h} = 13{,}333 \, \mathrm{m/s}\\
                & s_{\text{heen}} = 24 \cdot 10^{3} \, \mathrm{m} & & v_{\text{terug}} = 36 \, \mathrm{km/h} = 10 \, \mathrm{m/s}\\
                & s_{\text{terug}} = 6{,}0 \cdot 10^{3} \, \mathrm{m} & & t_{\text{wacht}} = 15 \, \mathrm{min} = 900 \, \mathrm{s}
            \end{align*}

            \underline{Gevraagd:} $\langle v_x \rangle$\\

            \underline{Oplossing:}\\
            We berekenen de duur van elk deel van de rit:
            \begin{align*}
                \Delta t_{\text{heen}} &= \frac{24 \cdot 10^{3} \, \mathrm{m}}{13{,}333 \, \mathrm{m/s}} = 1800 \, \mathrm{s}\\
                \Delta t_{\text{terug}} &= \frac{6{,}0 \cdot 10^{3} \, \mathrm{m}}{10 \, \mathrm{m/s}} = 600 \, \mathrm{s}
            \end{align*}
            De totale duur is $\Delta t = 1800 \, \mathrm{s} + 900 \, \mathrm{s} + 600 \, \mathrm{s} = 3300 \, \mathrm{s}$. De gemiddelde snelheid is de verplaatsing gedeeld door de duur:
            \begin{align*}
                \langle v_x \rangle &= \frac{\Delta x}{\Delta t} = \frac{18 \cdot 10^{3} \, \mathrm{m}}{3300 \, \mathrm{s}}\\
                &= 5{,}45 \, \mathrm{m/s} = 5{,}45 \cdot 3{,}6 \, \mathrm{km/h} = 19{,}6 \, \mathrm{km/h}
            \end{align*}
            Merk op dat dit niet het gemiddelde is van $48 \, \mathrm{km/h}$ en $36 \, \mathrm{km/h}$: de wachttijd en de zin van de beweging tellen mee.
        \end{solution}
    \end{question}

    \begin{question}
        Een boordcomputer meet de afgelegde weg gedeeld door de tijd. Bereken die gemiddelde snelheidsgrootte over de volledige rit.

        \[ \frac{s}{\Delta t} = \answer[tolerance=0.01,onlineshowanswerbutton]{32.7} \, \mathrm{km/h} \]

        \begin{solution}\nl
            \underline{Gegevens:}
            \begin{align*}
                & s = 30 \cdot 10^{3} \, \mathrm{m} & & \Delta t = 3300 \, \mathrm{s}
            \end{align*}

            \underline{Gevraagd:} $\dfrac{s}{\Delta t}$\\

            \underline{Oplossing:}\\
            \begin{align*}
                \frac{s}{\Delta t} &= \frac{30 \cdot 10^{3} \, \mathrm{m}}{3300 \, \mathrm{s}} = 9{,}09 \, \mathrm{m/s}\\
                &= 9{,}09 \cdot 3{,}6 \, \mathrm{km/h} = 32{,}7 \, \mathrm{km/h}
            \end{align*}
            Deze waarde is groter dan $\langle v_x \rangle = 19{,}6 \, \mathrm{km/h}$, omdat de afgelegde weg ($30 \, \mathrm{km}$) groter is dan de verplaatsing ($18 \, \mathrm{km}$).
        \end{solution}
    \end{question}

    \begin{question}
        De koerier wil dat zijn gemiddelde snelheid over de volledige rit (inclusief wachttijd) $20 \, \mathrm{km/h}$ bedraagt. Met welke constante snelheid had hij dan terug moeten rijden? De heenrit en de wachttijd blijven ongewijzigd.

        \[ v_{\text{terug}} = \answer[tolerance=0.01,onlineshowanswerbutton]{40} \, \mathrm{km/h} \]

        \begin{hint}
            Hoe lang mag de volledige rit hoogstens duren om op $20 \, \mathrm{km/h}$ uit te komen? Hoeveel van die tijd is al opgebruikt voor het heen en het wachten?
        \end{hint}

        \begin{solution}\nl
            \underline{Gegevens:}
            \begin{align*}
                & \Delta x = 18 \cdot 10^{3} \, \mathrm{m} & & \langle v_x \rangle = 20 \, \mathrm{km/h} = 5{,}556 \, \mathrm{m/s}\\
                & \Delta t_{\text{heen}} = 1800 \, \mathrm{s} & & t_{\text{wacht}} = 900 \, \mathrm{s}\\
                & s_{\text{terug}} = 6{,}0 \cdot 10^{3} \, \mathrm{m} & &
            \end{align*}

            \underline{Gevraagd:} $v_{\text{terug}}$\\

            \underline{Oplossing:}\\
            De verplaatsing blijft $18 \cdot 10^{3} \, \mathrm{m}$. De totale duur die bij $\langle v_x \rangle = 5{,}556 \, \mathrm{m/s}$ hoort, is:
            \[
                \Delta t = \frac{\Delta x}{\langle v_x \rangle} = \frac{18 \cdot 10^{3} \, \mathrm{m}}{5{,}556 \, \mathrm{m/s}} = 3240 \, \mathrm{s}
            \]
            Daarvan is al $1800 \, \mathrm{s} + 900 \, \mathrm{s} = 2700 \, \mathrm{s}$ opgebruikt. Voor de terugrit blijft dus $3240 \, \mathrm{s} - 2700 \, \mathrm{s} = 540 \, \mathrm{s}$ over:
            \[
                v_{\text{terug}} = \frac{6{,}0 \cdot 10^{3} \, \mathrm{m}}{540 \, \mathrm{s}} = 11{,}1 \, \mathrm{m/s} = 40 \, \mathrm{km/h}
            \]
        \end{solution}
    \end{question}
\end{exercise}
%———————————————————————————————————————

%———————————————————————————————————————
% OEFENING 2: v(t)-grafiek met zinsomkering
\begin{exercise}
    Een radiografisch bestuurbaar wagentje rijdt over een rechte rail langs de $x$-as. Op $t = 0{,}0 \, \mathrm{s}$ staat het op $x = +15 \, \mathrm{m}$. De figuur toont de snelheid $v_x$ als functie van de tijd.

    \begin{center}
        \begin{tikzpicture}[x=0.5cm,y=0.3cm]
            \draw[gray!40, very thin] (0,-8) grid[xstep=2,ystep=4] (20,8);
            \draw[positie] (0,0) -- (21,0) node[right] {$t\,(\mathrm{s})$};
            \draw[positie] (0,-9) -- (0,9.5) node[above] {$v_x\,(\mathrm{m/s})$};
            \foreach \tx in {2,4,...,20} \draw (\tx,0.3) -- (\tx,-0.3) node[below=3pt, font=\small] {\tx};
            \foreach \vy in {-8,-4,4,8} \draw (0.2,\vy) -- (-0.2,\vy) node[left, font=\small] {$\vy$};
            \draw[line width=1pt, red!100!black] (0,0) -- (4,8) -- (10,8) -- (14,-8) -- (18,-8) -- (20,0);
        \end{tikzpicture}
    \end{center}

    \begin{question}
        Bereken de versnelling $a_x$ tussen $t = 10 \, \mathrm{s}$ en $t = 14 \, \mathrm{s}$.

        \[ a_x = \answer[tolerance=0.01,onlineshowanswerbutton]{-4.0} \, \mathrm{m/s^2} \]

        Op $t = 13 \, \mathrm{s}$ geldt: \wordChoice{
            \choice[correct]{de grootte van de snelheid neemt toe}
            \choice{de grootte van de snelheid neemt af}
        }

        \begin{solution}\nl
            \underline{Gegevens:}
            \begin{align*}
                & v_x(10 \, \mathrm{s}) = 8{,}0 \, \mathrm{m/s} & & v_x(14 \, \mathrm{s}) = -8{,}0 \, \mathrm{m/s}
            \end{align*}

            \underline{Gevraagd:} $a_x$\\

            \underline{Oplossing:}\\
            De grafiek is in dit interval een rechte, dus de versnelling is constant en gelijk aan de helling:
            \[
                a_x = \frac{\Delta v_x}{\Delta t} = \frac{-8{,}0 \, \mathrm{m/s} - 8{,}0 \, \mathrm{m/s}}{14 \, \mathrm{s} - 10 \, \mathrm{s}} = -4{,}0 \, \mathrm{m/s^2}
            \]
            Op $t = 13 \, \mathrm{s}$ is de snelheid negatief ($v_x < 0$) en de versnelling ook negatief ($a_x < 0$). Ze hebben hetzelfde teken, dus de snelheid wordt (in grootte) groter: het wagentje versnelt in de negatieve zin.
        \end{solution}
    \end{question}

    \begin{question}
        Bereken de verplaatsing $\Delta x$ en de afgelegde weg $s$ tussen $t = 0 \, \mathrm{s}$ en $t = 20 \, \mathrm{s}$.

        \[ \Delta x = \answer[tolerance=0.01,onlineshowanswerbutton]{24} \, \mathrm{m} \qquad \text{en} \qquad s = \answer[tolerance=0.01,onlineshowanswerbutton]{120} \, \mathrm{m} \]

        \begin{hint}
            Wat stelt de oppervlakte tussen de grafiek en de $t$-as voor? Wat gebeurt er met het teken van die oppervlakte wanneer de snelheid negatief is?
        \end{hint}
        \begin{hint}
            Splits de grafiek op in driehoeken en rechthoeken. Op welk tijdstip verandert de zin van de beweging?
        \end{hint}

        \begin{solution}\nl
            \underline{Gegevens:}\\
            Rechtlijnige beweging; de snelheid wordt afgelezen uit de $v_x(t)$-grafiek. De zin van de beweging keert om op $t = 12 \, \mathrm{s}$ (daar is $v_x = 0$).\\

            \underline{Gevraagd:} $\Delta x$ en $s$\\

            \underline{Oplossing:}\\
            De verplaatsing is de oppervlakte onder de grafiek, met een negatief teken voor het deel onder de $t$-as. We splitsen op in zes stukken:
            \begin{align*}
                A_1 &= \tfrac{1}{2} \cdot 4{,}0 \, \mathrm{s} \cdot 8{,}0 \, \mathrm{m/s} = +16 \, \mathrm{m} & & (0 \, \mathrm{s} \text{ tot } 4 \, \mathrm{s})\\
                A_2 &= 6{,}0 \, \mathrm{s} \cdot 8{,}0 \, \mathrm{m/s} = +48 \, \mathrm{m} & & (4 \, \mathrm{s} \text{ tot } 10 \, \mathrm{s})\\
                A_3 &= \tfrac{1}{2} \cdot 2{,}0 \, \mathrm{s} \cdot 8{,}0 \, \mathrm{m/s} = +8{,}0 \, \mathrm{m} & & (10 \, \mathrm{s} \text{ tot } 12 \, \mathrm{s})\\
                A_4 &= \tfrac{1}{2} \cdot 2{,}0 \, \mathrm{s} \cdot (-8{,}0 \, \mathrm{m/s}) = -8{,}0 \, \mathrm{m} & & (12 \, \mathrm{s} \text{ tot } 14 \, \mathrm{s})\\
                A_5 &= 4{,}0 \, \mathrm{s} \cdot (-8{,}0 \, \mathrm{m/s}) = -32 \, \mathrm{m} & & (14 \, \mathrm{s} \text{ tot } 18 \, \mathrm{s})\\
                A_6 &= \tfrac{1}{2} \cdot 2{,}0 \, \mathrm{s} \cdot (-8{,}0 \, \mathrm{m/s}) = -8{,}0 \, \mathrm{m} & & (18 \, \mathrm{s} \text{ tot } 20 \, \mathrm{s})
            \end{align*}
            De verplaatsing is de som:
            \[
                \Delta x = 16 \, \mathrm{m} + 48 \, \mathrm{m} + 8{,}0 \, \mathrm{m} - 8{,}0 \, \mathrm{m} - 32 \, \mathrm{m} - 8{,}0 \, \mathrm{m} = 24 \, \mathrm{m}
            \]
            Voor de afgelegde weg tellen we alle oppervlaktes als positief:
            \[
                s = 16 \, \mathrm{m} + 48 \, \mathrm{m} + 8{,}0 \, \mathrm{m} + 8{,}0 \, \mathrm{m} + 32 \, \mathrm{m} + 8{,}0 \, \mathrm{m} = 120 \, \mathrm{m}
            \]
        \end{solution}
    \end{question}

    \begin{question}
        Bereken de gemiddelde snelheid $\langle v_x \rangle$ tussen $t = 0 \, \mathrm{s}$ en $t = 20 \, \mathrm{s}$.

        \[ \langle v_x \rangle = \answer[tolerance=0.01,onlineshowanswerbutton]{1.2} \, \mathrm{m/s} \]

        \begin{solution}\nl
            \underline{Gegevens:}
            \begin{align*}
                & \Delta x = 24 \, \mathrm{m} & & \Delta t = 20 \, \mathrm{s} - 0 \, \mathrm{s} = 20 \, \mathrm{s}
            \end{align*}

            \underline{Gevraagd:} $\langle v_x \rangle$\\

            \underline{Oplossing:}\\
            \[
                \langle v_x \rangle = \frac{\Delta x}{\Delta t} = \frac{24 \, \mathrm{m}}{20 \, \mathrm{s}} = 1{,}2 \, \mathrm{m/s}
            \]
            Het is dus de verplaatsing (en niet de afgelegde weg) die in de gemiddelde snelheid terechtkomt.
        \end{solution}
    \end{question}

    \begin{question}
        Wat is de grootste positie $x_{\max}$ die het wagentje tijdens de rit bereikt?

        \[ x_{\max} = \answer[tolerance=0.01,onlineshowanswerbutton]{87} \, \mathrm{m} \]

        \begin{hint}
            Op welk tijdstip stopt het wagentje met vooruitrijden? Vergeet de beginpositie niet.
        \end{hint}

        \begin{solution}\nl
            \underline{Gegevens:}
            \begin{align*}
                & x(0) = 15 \, \mathrm{m} & & A_1 + A_2 + A_3 = 72 \, \mathrm{m}
            \end{align*}

            \underline{Gevraagd:} $x_{\max}$\\

            \underline{Oplossing:}\\
            De positie is maximaal op het moment waarop de snelheid van teken wisselt van positief naar negatief: $t = 12 \, \mathrm{s}$. Tot dan is het wagentje $16 \, \mathrm{m} + 48 \, \mathrm{m} + 8{,}0 \, \mathrm{m} = 72 \, \mathrm{m}$ vooruitgereden, startend vanaf $15 \, \mathrm{m}$:
            \[
                x_{\max} = 15 \, \mathrm{m} + 72 \, \mathrm{m} = 87 \, \mathrm{m}
            \]
        \end{solution}
    \end{question}

    \begin{question}
        Het wagentje passeert de positie $x = 55 \, \mathrm{m}$ twee keer. Op welke tijdstippen $t_1$ en $t_2$ (met $t_1 < t_2$)?

        \[ t_1 = \answer[tolerance=0.01,onlineshowanswerbutton]{7.0} \, \mathrm{s} \qquad \text{en} \qquad t_2 = \answer[tolerance=0.01,onlineshowanswerbutton]{17} \, \mathrm{s} \]

        \begin{hint}
            Hoeveel moet het wagentje verplaatst zijn ten opzichte van zijn beginpositie om op $55 \, \mathrm{m}$ te staan? In welke tijdsintervallen kan dat gebeuren?
        \end{hint}

        \begin{solution}\nl
            \underline{Gegevens:}
            \begin{align*}
                & x(0) = 15 \, \mathrm{m} & & x = 55 \, \mathrm{m}
            \end{align*}

            \underline{Gevraagd:} $t_1$ en $t_2$\\

            \underline{Oplossing:}\\
            Het wagentje moet $55 \, \mathrm{m} - 15 \, \mathrm{m} = 40 \, \mathrm{m}$ verplaatst zijn ten opzichte van $t = 0 \, \mathrm{s}$.

            De verplaatsing na $4{,}0 \, \mathrm{s}$ is $16 \, \mathrm{m}$ en na $10 \, \mathrm{s}$ is dat $64 \, \mathrm{m}$. De eerste keer gebeurt dit dus tijdens de constante snelheid van $8{,}0 \, \mathrm{m/s}$:
            \[
                16 \, \mathrm{m} + 8{,}0 \, \mathrm{m/s} \cdot (t_1 - 4{,}0 \, \mathrm{s}) = 40 \, \mathrm{m} \quad \Rightarrow \quad t_1 = 4{,}0 \, \mathrm{s} + \frac{24 \, \mathrm{m}}{8{,}0 \, \mathrm{m/s}} = 7{,}0 \, \mathrm{s}
            \]
            Na $t = 12 \, \mathrm{s}$ rijdt het wagentje terug. Op $t = 14 \, \mathrm{s}$ is de verplaatsing nog $72 \, \mathrm{m} - 8{,}0 \, \mathrm{m} = 64 \, \mathrm{m}$, en tot $t = 18 \, \mathrm{s}$ rijdt het met constante snelheid $-8{,}0 \, \mathrm{m/s}$ terug:
            \[
                64 \, \mathrm{m} + (-8{,}0 \, \mathrm{m/s}) \cdot (t_2 - 14 \, \mathrm{s}) = 40 \, \mathrm{m} \quad \Rightarrow \quad t_2 = 14 \, \mathrm{s} + \frac{24 \, \mathrm{m}}{8{,}0 \, \mathrm{m/s}} = 17 \, \mathrm{s}
            \]
        \end{solution}
    \end{question}
\end{exercise}
%———————————————————————————————————————

%———————————————————————————————————————
% OEFENING 3: a(t)-grafiek, van versnelling naar snelheid naar positie
\begin{exercise}
    Een voorwerp beweegt langs de $x$-as. Op $t = 0{,}0 \, \mathrm{s}$ bevindt het zich in de oorsprong en heeft het een snelheid $v_x = -6{,}0 \, \mathrm{m/s}$. De figuur toont de versnelling $a_x$ als functie van de tijd.

    \begin{center}
        \begin{tikzpicture}[x=0.6cm,y=0.5cm]
            \draw[gray!40, very thin] (0,-2) grid[xstep=1,ystep=1] (10,3);
            \draw[positie] (0,0) -- (11,0) node[right] {$t\,(\mathrm{s})$};
            \draw[positie] (0,-2.7) -- (0,3.8) node[above] {$a_x\,(\mathrm{m/s^2})$};
            \foreach \tx in {2,4,...,10} \draw (\tx,0.1) -- (\tx,-0.1) node[below=3pt, font=\small] {\tx};
            \foreach \ay in {-2,-1,1,2,3} \draw (0.1,\ay) -- (-0.1,\ay) node[left, font=\small] {$\ay$};
            \draw[line width=1pt, violet!100!black] (0,3) -- (4,3);
            \draw[line width=1pt, violet!100!black] (4,0) -- (6,0);
            \draw[line width=1pt, violet!100!black] (6,-2) -- (10,-2);
            \draw[dotted, violet!100!black] (4,3) -- (4,0) (6,0) -- (6,-2);
        \end{tikzpicture}
    \end{center}

    \begin{question}
        Bereken de snelheid $v_x$ op $t = 4{,}0 \, \mathrm{s}$.

        \[ v_x(4{,}0 \, \mathrm{s}) = \answer[tolerance=0.01,onlineshowanswerbutton]{6.0} \, \mathrm{m/s} \]

        \begin{hint}
            Wat stelt de oppervlakte onder een $a_x(t)$-grafiek voor?
        \end{hint}

        \begin{solution}\nl
            \underline{Gegevens:}
            \begin{align*}
                & v_x(0) = -6{,}0 \, \mathrm{m/s} & & a_x = 3{,}0 \, \mathrm{m/s^2} \text{ voor } 0 \, \mathrm{s} \leq t \leq 4{,}0 \, \mathrm{s}
            \end{align*}

            \underline{Gevraagd:} $v_x(4{,}0 \, \mathrm{s})$\\

            \underline{Oplossing:}\\
            De oppervlakte onder de $a_x(t)$-grafiek is de snelheidsverandering:
            \[
                \Delta v_x = 3{,}0 \, \mathrm{m/s^2} \cdot 4{,}0 \, \mathrm{s} = 12 \, \mathrm{m/s}
            \]
            \[
                v_x(4{,}0 \, \mathrm{s}) = v_x(0) + \Delta v_x = -6{,}0 \, \mathrm{m/s} + 12 \, \mathrm{m/s} = 6{,}0 \, \mathrm{m/s}
            \]
        \end{solution}
    \end{question}

    \begin{question}
        Het voorwerp keert tweemaal van zin om. Op welke tijdstippen $t_1$ en $t_2$ (met $t_1 < t_2$)?

        \[ t_1 = \answer[tolerance=0.01,onlineshowanswerbutton]{2.0} \, \mathrm{s} \qquad \text{en} \qquad t_2 = \answer[tolerance=0.01,onlineshowanswerbutton]{9.0} \, \mathrm{s} \]

        \begin{hint}
            Wat is de snelheid op het moment dat het voorwerp van zin verandert?
        \end{hint}
        \begin{hint}
            Welke snelheid heeft het voorwerp op $t = 6{,}0 \, \mathrm{s}$, en hoe verandert die daarna?
        \end{hint}

        \begin{solution}\nl
            \underline{Gegevens:}
            \begin{align*}
                & v_x(0) = -6{,}0 \, \mathrm{m/s} & & v_x(4{,}0 \, \mathrm{s}) = 6{,}0 \, \mathrm{m/s}
            \end{align*}

            \underline{Gevraagd:} $t_1$ en $t_2$\\

            \underline{Oplossing:}\\
            Bij een zinsomkering is $v_x = 0$. In het eerste interval ($a_x = 3{,}0 \, \mathrm{m/s^2}$) geldt:
            \[
                -6{,}0 \, \mathrm{m/s} + 3{,}0 \, \mathrm{m/s^2} \cdot t_1 = 0 \quad \Rightarrow \quad t_1 = \frac{6{,}0 \, \mathrm{m/s}}{3{,}0 \, \mathrm{m/s^2}} = 2{,}0 \, \mathrm{s}
            \]
            Tussen $4{,}0 \, \mathrm{s}$ en $6{,}0 \, \mathrm{s}$ is $a_x = 0$, dus blijft de snelheid $6{,}0 \, \mathrm{m/s}$. Vanaf $t = 6{,}0 \, \mathrm{s}$ is $a_x = -2{,}0 \, \mathrm{m/s^2}$:
            \[
                6{,}0 \, \mathrm{m/s} - 2{,}0 \, \mathrm{m/s^2} \cdot (t_2 - 6{,}0 \, \mathrm{s}) = 0 \quad \Rightarrow \quad t_2 = 6{,}0 \, \mathrm{s} + \frac{6{,}0 \, \mathrm{m/s}}{2{,}0 \, \mathrm{m/s^2}} = 9{,}0 \, \mathrm{s}
            \]
        \end{solution}
    \end{question}

    \begin{question}
        Bereken de verplaatsing $\Delta x$ en de afgelegde weg $s$ tussen $t = 0 \, \mathrm{s}$ en $t = 10 \, \mathrm{s}$.

        \[ \Delta x = \answer[tolerance=0.01,onlineshowanswerbutton]{20} \, \mathrm{m} \qquad \text{en} \qquad s = \answer[tolerance=0.01,onlineshowanswerbutton]{34} \, \mathrm{m} \]

        \begin{hint}
            Schets eerst de $v_x(t)$-grafiek. Gebruik de snelheden op $t = 0$, $2{,}0$, $4{,}0$, $6{,}0$, $9{,}0$ en $10 \, \mathrm{s}$.
        \end{hint}
        \begin{hint}
            De oppervlakte onder de $v_x(t)$-grafiek is de verplaatsing. Rekening houden met het teken geeft de verplaatsing, alles positief tellen geeft de afgelegde weg.
        \end{hint}

        \begin{solution}\nl
            \underline{Gegevens:}
            \begin{align*}
                & v_x(0) = -6{,}0 \, \mathrm{m/s} & & v_x(2{,}0 \, \mathrm{s}) = 0 \, \mathrm{m/s}\\
                & v_x(4{,}0 \, \mathrm{s}) = v_x(6{,}0 \, \mathrm{s}) = 6{,}0 \, \mathrm{m/s} & & v_x(9{,}0 \, \mathrm{s}) = 0 \, \mathrm{m/s}\\
                & v_x(10 \, \mathrm{s}) = -2{,}0 \, \mathrm{m/s} & &
            \end{align*}

            \underline{Gevraagd:} $\Delta x$ en $s$\\

            \underline{Oplossing:}\\
            De $v_x(t)$-grafiek bestaat uit rechte stukken. De oppervlaktes zijn:
            \begin{align*}
                A_1 &= \tfrac{1}{2} \cdot 2{,}0 \, \mathrm{s} \cdot (-6{,}0 \, \mathrm{m/s}) = -6{,}0 \, \mathrm{m} & & (0 \, \mathrm{s} \text{ tot } 2{,}0 \, \mathrm{s})\\
                A_2 &= \tfrac{1}{2} \cdot 2{,}0 \, \mathrm{s} \cdot 6{,}0 \, \mathrm{m/s} = +6{,}0 \, \mathrm{m} & & (2{,}0 \, \mathrm{s} \text{ tot } 4{,}0 \, \mathrm{s})\\
                A_3 &= 2{,}0 \, \mathrm{s} \cdot 6{,}0 \, \mathrm{m/s} = +12 \, \mathrm{m} & & (4{,}0 \, \mathrm{s} \text{ tot } 6{,}0 \, \mathrm{s})\\
                A_4 &= \tfrac{1}{2} \cdot 3{,}0 \, \mathrm{s} \cdot 6{,}0 \, \mathrm{m/s} = +9{,}0 \, \mathrm{m} & & (6{,}0 \, \mathrm{s} \text{ tot } 9{,}0 \, \mathrm{s})\\
                A_5 &= \tfrac{1}{2} \cdot 1{,}0 \, \mathrm{s} \cdot (-2{,}0 \, \mathrm{m/s}) = -1{,}0 \, \mathrm{m} & & (9{,}0 \, \mathrm{s} \text{ tot } 10 \, \mathrm{s})
            \end{align*}
            De verplaatsing is de som met tekens:
            \[
                \Delta x = -6{,}0 \, \mathrm{m} + 6{,}0 \, \mathrm{m} + 12 \, \mathrm{m} + 9{,}0 \, \mathrm{m} - 1{,}0 \, \mathrm{m} = 20 \, \mathrm{m}
            \]
            Voor de afgelegde weg tellen we alles positief:
            \[
                s = 6{,}0 \, \mathrm{m} + 6{,}0 \, \mathrm{m} + 12 \, \mathrm{m} + 9{,}0 \, \mathrm{m} + 1{,}0 \, \mathrm{m} = 34 \, \mathrm{m}
            \]
        \end{solution}
    \end{question}

    \begin{question}
        Wat is de grootste afstand tussen het voorwerp en de oorsprong tijdens deze $10 \, \mathrm{s}$?

        \[ \answer[tolerance=0.01,onlineshowanswerbutton]{21} \, \mathrm{m} \]

        \begin{hint}
            Bepaal de positie op elk tijdstip waarop het voorwerp van zin keert. Welke van die posities ligt het verst van de oorsprong?
        \end{hint}

        \begin{solution}\nl
            \underline{Gegevens:}
            \begin{align*}
                & x(0) = 0 \, \mathrm{m} & & t_1 = 2{,}0 \, \mathrm{s}, \quad t_2 = 9{,}0 \, \mathrm{s}
            \end{align*}

            \underline{Gevraagd:} grootste afstand tot de oorsprong\\

            \underline{Oplossing:}\\
            De positie is extreem op de momenten van zinsomkering:
            \begin{align*}
                x(2{,}0 \, \mathrm{s}) &= A_1 = -6{,}0 \, \mathrm{m}\\
                x(9{,}0 \, \mathrm{s}) &= A_1 + A_2 + A_3 + A_4 = -6{,}0 \, \mathrm{m} + 6{,}0 \, \mathrm{m} + 12 \, \mathrm{m} + 9{,}0 \, \mathrm{m} = 21 \, \mathrm{m}
            \end{align*}
            Op $t = 10 \, \mathrm{s}$ is $x = 20 \, \mathrm{m}$. De grootste afstand tot de oorsprong is dus $21 \, \mathrm{m}$ (op $t = 9{,}0 \, \mathrm{s}$), en niet de eindpositie.
        \end{solution}
    \end{question}
\end{exercise}
%———————————————————————————————————————

%———————————————————————————————————————
% OEFENING 4: positievergelijking (derde graad), afgeleiden en zinsomkering
\begin{exercise}
    De positie van een voorwerp op de $x$-as wordt gegeven door
    \[
        x(t) = 2{,}0 \, \mathrm{\frac{m}{s^3}} \cdot t^{3} - 9{,}0 \, \mathrm{\frac{m}{s^2}} \cdot t^{2} + 12 \, \mathrm{\frac{m}{s}} \cdot t + 1{,}0 \, \mathrm{m}
    \]
    met $t$ in seconden. We bekijken de beweging tussen $t = 0 \, \mathrm{s}$ en $t = 3{,}0 \, \mathrm{s}$.

    \begin{question}
        Het voorwerp staat twee keer even stil. Op welke tijdstippen $t_1$ en $t_2$ (met $t_1 < t_2$)?

        \[ t_1 = \answer[tolerance=0.01,onlineshowanswerbutton]{1.0} \, \mathrm{s} \qquad \text{en} \qquad t_2 = \answer[tolerance=0.01,onlineshowanswerbutton]{2.0} \, \mathrm{s} \]

        \begin{hint}
            Welke grootheid is nul als een voorwerp even stilstaat, en hoe vind je de vergelijking van die grootheid vanuit $x(t)$?
        \end{hint}

        \begin{solution}\nl
            \underline{Gegevens:}\\
            $x(t) = 2{,}0 \, \mathrm{m/s^3} \cdot t^{3} - 9{,}0 \, \mathrm{m/s^2} \cdot t^{2} + 12 \, \mathrm{m/s} \cdot t + 1{,}0 \, \mathrm{m}$\\

            \underline{Gevraagd:} de tijdstippen waarvoor $v_x = 0$\\

            \underline{Oplossing:}\\
            De snelheid is de afgeleide van de positie:
            \[
                v_x(t) = x'(t) = 6{,}0 \, \mathrm{m/s^3} \cdot t^{2} - 18 \, \mathrm{m/s^2} \cdot t + 12 \, \mathrm{m/s}
            \]
            Stilstaan betekent $v_x = 0$. We laten de eenheden even weg ($t$ in $\mathrm{s}$) en delen door $6$:
            \[
                t^{2} - 3t + 2 = 0 \quad \Rightarrow \quad t = \frac{3 \pm \sqrt{9 - 8}}{2} \quad \Rightarrow \quad t_1 = 1{,}0 \, \mathrm{s} \text{ en } t_2 = 2{,}0 \, \mathrm{s}
            \]
        \end{solution}
    \end{question}

    \begin{question}
        Bereken de verplaatsing $\Delta x$ tussen $t = 0 \, \mathrm{s}$ en $t = 3{,}0 \, \mathrm{s}$.

        \[ \Delta x = \answer[tolerance=0.01,onlineshowanswerbutton]{9.0} \, \mathrm{m} \]

        \begin{solution}\nl
            \underline{Gegevens:}\\
            $x(t) = 2{,}0 \, \mathrm{m/s^3} \cdot t^{3} - 9{,}0 \, \mathrm{m/s^2} \cdot t^{2} + 12 \, \mathrm{m/s} \cdot t + 1{,}0 \, \mathrm{m}$\\

            \underline{Gevraagd:} $\Delta x$\\

            \underline{Oplossing:}\\
            We berekenen de positie op beide tijdstippen:
            \begin{align*}
                x(0) &= 1{,}0 \, \mathrm{m}\\
                x(3{,}0 \, \mathrm{s}) &= 2{,}0 \, \mathrm{m/s^3} \cdot (3{,}0 \, \mathrm{s})^{3} - 9{,}0 \, \mathrm{m/s^2} \cdot (3{,}0 \, \mathrm{s})^{2} + 12 \, \mathrm{m/s} \cdot 3{,}0 \, \mathrm{s} + 1{,}0 \, \mathrm{m}\\
                &= 54 \, \mathrm{m} - 81 \, \mathrm{m} + 36 \, \mathrm{m} + 1{,}0 \, \mathrm{m} = 10 \, \mathrm{m}
            \end{align*}
            \[
                \Delta x = x(3{,}0 \, \mathrm{s}) - x(0) = 10 \, \mathrm{m} - 1{,}0 \, \mathrm{m} = 9{,}0 \, \mathrm{m}
            \]
        \end{solution}
    \end{question}

    \begin{question}
        Bereken de afgelegde weg $s$ tussen $t = 0 \, \mathrm{s}$ en $t = 3{,}0 \, \mathrm{s}$.

        \[ s = \answer[tolerance=0.01,onlineshowanswerbutton]{11} \, \mathrm{m} \]

        \begin{hint}
            De afgelegde weg is niet gelijk aan de verplaatsing als het voorwerp van zin verandert. Splits het tijdsinterval op in stukken waarin het voorwerp telkens in dezelfde zin beweegt.
        \end{hint}

        \begin{solution}\nl
            \underline{Gegevens:}\\
            Zinsomkeringen op $t_1 = 1{,}0 \, \mathrm{s}$ en $t_2 = 2{,}0 \, \mathrm{s}$; $x(0) = 1{,}0 \, \mathrm{m}$ en $x(3{,}0 \, \mathrm{s}) = 10 \, \mathrm{m}$.\\

            \underline{Gevraagd:} $s$\\

            \underline{Oplossing:}\\
            We berekenen de posities op de momenten van zinsomkering:
            \begin{align*}
                x(1{,}0 \, \mathrm{s}) &= 2{,}0 \, \mathrm{m} - 9{,}0 \, \mathrm{m} + 12 \, \mathrm{m} + 1{,}0 \, \mathrm{m} = 6{,}0 \, \mathrm{m}\\
                x(2{,}0 \, \mathrm{s}) &= 16 \, \mathrm{m} - 36 \, \mathrm{m} + 24 \, \mathrm{m} + 1{,}0 \, \mathrm{m} = 5{,}0 \, \mathrm{m}
            \end{align*}
            Het voorwerp gaat dus van $1{,}0 \, \mathrm{m}$ naar $6{,}0 \, \mathrm{m}$, dan terug naar $5{,}0 \, \mathrm{m}$ en daarna naar $10 \, \mathrm{m}$:
            \[
                s = \lvert 6{,}0 \, \mathrm{m} - 1{,}0 \, \mathrm{m} \rvert + \lvert 5{,}0 \, \mathrm{m} - 6{,}0 \, \mathrm{m} \rvert + \lvert 10 \, \mathrm{m} - 5{,}0 \, \mathrm{m} \rvert = 5{,}0 \, \mathrm{m} + 1{,}0 \, \mathrm{m} + 5{,}0 \, \mathrm{m} = 11 \, \mathrm{m}
            \]
        \end{solution}
    \end{question}

    \begin{question}
        Wat is de snelheid $v_x$ van het voorwerp op het moment dat zijn versnelling nul is?

        \[ v_x = \answer[tolerance=0.01,onlineshowanswerbutton]{-1.5} \, \mathrm{m/s} \]

        \begin{hint}
            Hoe vind je de versnellingsvergelijking $a_x(t)$ vanuit $x(t)$? Op welk tijdstip is die nul?
        \end{hint}

        \begin{solution}\nl
            \underline{Gegevens:}\\
            $v_x(t) = 6{,}0 \, \mathrm{m/s^3} \cdot t^{2} - 18 \, \mathrm{m/s^2} \cdot t + 12 \, \mathrm{m/s}$\\

            \underline{Gevraagd:} $v_x$ als $a_x = 0$\\

            \underline{Oplossing:}\\
            De versnelling is de afgeleide van de snelheid:
            \[
                a_x(t) = v_x'(t) = 12 \, \mathrm{m/s^3} \cdot t - 18 \, \mathrm{m/s^2}
            \]
            Dit is nul voor $t = \dfrac{18 \, \mathrm{m/s^2}}{12 \, \mathrm{m/s^3}} = 1{,}5 \, \mathrm{s}$. De snelheid op dat moment is:
            \begin{align*}
                v_x(1{,}5 \, \mathrm{s}) &= 6{,}0 \, \mathrm{m/s^3} \cdot (1{,}5 \, \mathrm{s})^{2} - 18 \, \mathrm{m/s^2} \cdot 1{,}5 \, \mathrm{s} + 12 \, \mathrm{m/s}\\
                &= 13{,}5 \, \mathrm{m/s} - 27 \, \mathrm{m/s} + 12 \, \mathrm{m/s} = -1{,}5 \, \mathrm{m/s}
            \end{align*}
            Het voorwerp beweegt op dat moment in de negatieve zin, en zijn snelheid is op dat moment het kleinst (meest negatief).
        \end{solution}
    \end{question}

    \begin{question}
        Op twee tijdstippen tussen $0 \, \mathrm{s}$ en $3{,}0 \, \mathrm{s}$ is de ogenblikkelijke snelheid exact gelijk aan de gemiddelde snelheid over dat interval. Bepaal die tijdstippen $t_1$ en $t_2$ (met $t_1 < t_2$).

        \[ t_1 = \answer[tolerance=0.02,onlineshowanswerbutton]{0.63} \, \mathrm{s} \qquad \text{en} \qquad t_2 = \answer[tolerance=0.02,onlineshowanswerbutton]{2.37} \, \mathrm{s} \]

        \begin{hint}
            Bereken eerst de gemiddelde snelheid over het interval. Stel daarna de snelheidsvergelijking gelijk aan die waarde.
        \end{hint}

        \begin{solution}\nl
            \underline{Gegevens:}
            \begin{align*}
                & \Delta x = 9{,}0 \, \mathrm{m} & & \Delta t = 3{,}0 \, \mathrm{s}
            \end{align*}

            \underline{Gevraagd:} $t_1$ en $t_2$ waarvoor $v_x(t) = \langle v_x \rangle$\\

            \underline{Oplossing:}\\
            De gemiddelde snelheid is:
            \[
                \langle v_x \rangle = \frac{\Delta x}{\Delta t} = \frac{9{,}0 \, \mathrm{m}}{3{,}0 \, \mathrm{s}} = 3{,}0 \, \mathrm{m/s}
            \]
            We stellen $v_x(t) = 3{,}0 \, \mathrm{m/s}$ (eenheden even weggelaten, $t$ in $\mathrm{s}$):
            \[
                6t^{2} - 18t + 12 = 3 \quad \Rightarrow \quad 2t^{2} - 6t + 3 = 0
            \]
            \[
                t = \frac{6 \pm \sqrt{36 - 24}}{4} = \frac{6 \pm 3{,}464}{4} \quad \Rightarrow \quad t_1 = 0{,}63 \, \mathrm{s} \text{ en } t_2 = 2{,}37 \, \mathrm{s}
            \]
            Beide tijdstippen liggen binnen het interval $[0 \, \mathrm{s}; 3{,}0 \, \mathrm{s}]$.
        \end{solution}
    \end{question}
\end{exercise}
%———————————————————————————————————————

%———————————————————————————————————————
% OEFENING 5: onbekende coëfficiënten bepalen uit voorwaarden
\begin{exercise}
    De positie van een voorwerp op de $x$-as wordt gegeven door
    \[
        x(t) = \alpha \cdot t^{3} + \beta \cdot t^{2}
    \]
    waarbij $\alpha$ en $\beta$ constanten zijn. Op $t = 0 \, \mathrm{s}$ bevindt het voorwerp zich in de oorsprong. Op $t = 2{,}0 \, \mathrm{s}$ staat het even stil op de positie $x = -8{,}0 \, \mathrm{m}$.

    \begin{question}
        Bepaal $\alpha$ en $\beta$.

        \[ \alpha = \answer[tolerance=0.01,onlineshowanswerbutton]{2.0} \, \mathrm{m/s^3} \qquad \text{en} \qquad \beta = \answer[tolerance=0.01,onlineshowanswerbutton]{-6.0} \, \mathrm{m/s^2} \]

        \begin{hint}
            Je hebt twee onbekenden, dus twee voorwaarden nodig. Welke twee gegevens uit de opgave kan je in wiskundige vorm schrijven?
        \end{hint}
        \begin{hint}
            Voor een van de voorwaarden heb je de snelheidsvergelijking nodig. Hoe vind je die?
        \end{hint}

        \begin{solution}\nl
            \underline{Gegevens:}
            \begin{align*}
                & x(2{,}0 \, \mathrm{s}) = -8{,}0 \, \mathrm{m} & & v_x(2{,}0 \, \mathrm{s}) = 0 \, \mathrm{m/s}
            \end{align*}

            \underline{Gevraagd:} $\alpha$ en $\beta$\\

            \underline{Oplossing:}\\
            De snelheid is de afgeleide van de positie:
            \[
                v_x(t) = 3\alpha \cdot t^{2} + 2\beta \cdot t
            \]
            Voorwaarde 1 (even stilstaan op $t = 2{,}0 \, \mathrm{s}$), met $t$ in $\mathrm{s}$:
            \[
                v_x(2{,}0) = 12\alpha + 4\beta = 0 \quad \Rightarrow \quad \beta = -3\alpha
            \]
            Voorwaarde 2 (positie op $t = 2{,}0 \, \mathrm{s}$):
            \[
                x(2{,}0) = 8\alpha + 4\beta = -8{,}0
            \]
            We vullen $\beta = -3\alpha$ in: $8\alpha - 12\alpha = -4\alpha = -8{,}0$, dus $\alpha = 2{,}0$ en $\beta = -6{,}0$.

            Uit de eenheid van $x$ (in $\mathrm{m}$) volgt: $[\alpha] = \mathrm{m/s^3}$ en $[\beta] = \mathrm{m/s^2}$. Dus
            \[
                x(t) = 2{,}0 \, \mathrm{\frac{m}{s^3}} \cdot t^{3} - 6{,}0 \, \mathrm{\frac{m}{s^2}} \cdot t^{2}
            \]
        \end{solution}
    \end{question}

    \begin{question}
        Na hoeveel tijd keert het voorwerp (na het vertrek) terug in de oorsprong?

        \[ t = \answer[tolerance=0.01,onlineshowanswerbutton]{3.0} \, \mathrm{s} \]

        \begin{hint}
            Je kan de positievergelijking ontbinden in factoren.
        \end{hint}

        \begin{solution}\nl
            \underline{Gegevens:}\\
            $x(t) = 2{,}0 \, \mathrm{m/s^3} \cdot t^{3} - 6{,}0 \, \mathrm{m/s^2} \cdot t^{2}$\\

            \underline{Gevraagd:} $t$ waarvoor $x = 0$ (met $t > 0$)\\

            \underline{Oplossing:}\\
            We stellen $x(t) = 0$ (eenheden even weggelaten) en zetten $2t^{2}$ buiten haakjes:
            \[
                2t^{3} - 6t^{2} = 2t^{2} \cdot (t - 3) = 0
            \]
            De oplossing $t = 0$ is het vertrek zelf. Het voorwerp is dus terug in de oorsprong op $t = 3{,}0 \, \mathrm{s}$.
        \end{solution}
    \end{question}

    \begin{question}
        Bereken de snelheid $v_x$ waarmee het voorwerp terug door de oorsprong gaat.

        \[ v_x = \answer[tolerance=0.01,onlineshowanswerbutton]{18} \, \mathrm{m/s} \]

        \begin{solution}\nl
            \underline{Gegevens:}\\
            $v_x(t) = 6{,}0 \, \mathrm{m/s^3} \cdot t^{2} - 12 \, \mathrm{m/s^2} \cdot t$ en $t = 3{,}0 \, \mathrm{s}$\\

            \underline{Gevraagd:} $v_x(3{,}0 \, \mathrm{s})$\\

            \underline{Oplossing:}\\
            \begin{align*}
                v_x(3{,}0 \, \mathrm{s}) &= 6{,}0 \, \mathrm{m/s^3} \cdot (3{,}0 \, \mathrm{s})^{2} - 12 \, \mathrm{m/s^2} \cdot 3{,}0 \, \mathrm{s}\\
                &= 54 \, \mathrm{m/s} - 36 \, \mathrm{m/s} = 18 \, \mathrm{m/s}
            \end{align*}
            Het voorwerp gaat in de positieve zin door de oorsprong.
        \end{solution}
    \end{question}

    \begin{question}
        Bereken de afgelegde weg $s$ tussen $t = 0 \, \mathrm{s}$ en het moment waarop het voorwerp terug in de oorsprong is.

        \[ s = \answer[tolerance=0.01,onlineshowanswerbutton]{16} \, \mathrm{m} \]

        \begin{hint}
            De verplaatsing over dit interval is nul. Is de afgelegde weg dat dan ook? Bepaal wanneer het voorwerp van zin keert.
        \end{hint}

        \begin{solution}\nl
            \underline{Gegevens:}\\
            $v_x(t) = 6{,}0 \, \mathrm{m/s^3} \cdot t^{2} - 12 \, \mathrm{m/s^2} \cdot t = 6{,}0 \, \mathrm{m/s^3} \cdot t \cdot (t - 2{,}0 \, \mathrm{s})$; $x(2{,}0 \, \mathrm{s}) = -8{,}0 \, \mathrm{m}$\\

            \underline{Gevraagd:} $s$\\

            \underline{Oplossing:}\\
            De snelheid is nul op $t = 0 \, \mathrm{s}$ en $t = 2{,}0 \, \mathrm{s}$. Voor $0 \, \mathrm{s} < t < 2{,}0 \, \mathrm{s}$ is $v_x < 0$ en voor $t > 2{,}0 \, \mathrm{s}$ is $v_x > 0$. Het voorwerp beweegt dus eerst van $0 \, \mathrm{m}$ naar $-8{,}0 \, \mathrm{m}$ en keert dan terug naar $0 \, \mathrm{m}$:
            \[
                s = \lvert -8{,}0 \, \mathrm{m} - 0 \, \mathrm{m} \rvert + \lvert 0 \, \mathrm{m} - (-8{,}0 \, \mathrm{m}) \rvert = 8{,}0 \, \mathrm{m} + 8{,}0 \, \mathrm{m} = 16 \, \mathrm{m}
            \]
            Hoewel de verplaatsing en de gemiddelde snelheid nul zijn, is de afgelegde weg $16 \, \mathrm{m}$.
        \end{solution}
    \end{question}
\end{exercise}
%———————————————————————————————————————

%———————————————————————————————————————
% OEFENING 6: twee voertuigen, twee ontmoetingen (kwadratische positievergelijkingen)
\begin{exercise}
    Een motorrijder en een vrachtwagen rijden in dezelfde zin op dezelfde rechte weg (dezelfde $x$-as). Op $t = 0 \, \mathrm{s}$ bevindt de motor zich in de oorsprong en de vrachtwagen $27 \, \mathrm{m}$ verder. Hun posities worden gegeven door
    \begin{align*}
        x_M(t) &= 20 \, \mathrm{\frac{m}{s}} \cdot t - 1{,}0 \, \mathrm{\frac{m}{s^2}} \cdot t^{2}\\
        x_V(t) &= 27 \, \mathrm{m} + 8{,}0 \, \mathrm{\frac{m}{s}} \cdot t
    \end{align*}
    met $t$ in seconden. Deze vergelijkingen gelden voor $0 \, \mathrm{s} \leq t \leq 10 \, \mathrm{s}$.

    \begin{question}
        Wat is de beginsnelheid van de motor in $\mathrm{km/h}$?

        \[ v_{M}(0) = \answer[tolerance=0.01,onlineshowanswerbutton]{72} \, \mathrm{km/h} \]

        \begin{solution}\nl
            \underline{Gegevens:}\\
            $x_M(t) = 20 \, \mathrm{m/s} \cdot t - 1{,}0 \, \mathrm{m/s^2} \cdot t^{2}$\\

            \underline{Gevraagd:} $v_M(0)$ in $\mathrm{km/h}$\\

            \underline{Oplossing:}\\
            De snelheid is de afgeleide van de positie:
            \[
                v_M(t) = x_M'(t) = 20 \, \mathrm{m/s} - 2{,}0 \, \mathrm{m/s^2} \cdot t
            \]
            Dus $v_M(0) = 20 \, \mathrm{m/s} = 20 \cdot 3{,}6 \, \mathrm{km/h} = 72 \, \mathrm{km/h}$.
        \end{solution}
    \end{question}

    \begin{question}
        Op welke tijdstippen $t_1$ en $t_2$ (met $t_1 < t_2$) bevinden de motor en de vrachtwagen zich op dezelfde plaats?

        \[ t_1 = \answer[tolerance=0.01,onlineshowanswerbutton]{3.0} \, \mathrm{s} \qquad \text{en} \qquad t_2 = \answer[tolerance=0.01,onlineshowanswerbutton]{9.0} \, \mathrm{s} \]

        \begin{hint}
            Wat geldt er voor de posities van beide voertuigen op het moment van een ontmoeting? Dit levert een vergelijking in $t$ op die je moet oplossen.
        \end{hint}

        \begin{solution}\nl
            \underline{Gegevens:}\\
            $x_M(t) = 20 \, \mathrm{m/s} \cdot t - 1{,}0 \, \mathrm{m/s^2} \cdot t^{2}$ en $x_V(t) = 27 \, \mathrm{m} + 8{,}0 \, \mathrm{m/s} \cdot t$\\

            \underline{Gevraagd:} $t$ waarvoor $x_M = x_V$\\

            \underline{Oplossing:}\\
            We stellen de posities gelijk (eenheden even weggelaten, $t$ in $\mathrm{s}$):
            \[
                20t - t^{2} = 27 + 8t \quad \Rightarrow \quad t^{2} - 12t + 27 = 0
            \]
            \[
                t = \frac{12 \pm \sqrt{144 - 108}}{2} = \frac{12 \pm 6}{2} \quad \Rightarrow \quad t_1 = 3{,}0 \, \mathrm{s} \text{ en } t_2 = 9{,}0 \, \mathrm{s}
            \]
        \end{solution}
    \end{question}

    \begin{question}
        Bereken het snelheidsverschil $v_M - v_V$ op het eerste ontmoetingsmoment $t_1$.

        \[ v_M(t_1) - v_V(t_1) = \answer[tolerance=0.01,onlineshowanswerbutton]{6.0} \, \mathrm{m/s} \]

        \begin{solution}\nl
            \underline{Gegevens:}\\
            $v_M(t) = 20 \, \mathrm{m/s} - 2{,}0 \, \mathrm{m/s^2} \cdot t$ en $v_V(t) = 8{,}0 \, \mathrm{m/s}$; $t_1 = 3{,}0 \, \mathrm{s}$\\

            \underline{Gevraagd:} $v_M(t_1) - v_V(t_1)$\\

            \underline{Oplossing:}\\
            \begin{align*}
                v_M(3{,}0 \, \mathrm{s}) &= 20 \, \mathrm{m/s} - 2{,}0 \, \mathrm{m/s^2} \cdot 3{,}0 \, \mathrm{s} = 14 \, \mathrm{m/s}\\
                v_M - v_V &= 14 \, \mathrm{m/s} - 8{,}0 \, \mathrm{m/s} = 6{,}0 \, \mathrm{m/s}
            \end{align*}
            De motor is op dat moment sneller dan de vrachtwagen en haalt hem in.
        \end{solution}
    \end{question}

    \begin{question}
        Wat gebeurt er op het tweede ontmoetingsmoment $t_2$?

        \wordChoice{
            \choice{De motor haalt de vrachtwagen in.}
            \choice[correct]{De vrachtwagen haalt de motor in.}
            \choice{De twee voertuigen kruisen elkaar in tegengestelde zin.}
        }

        \begin{hint}
            Vergelijk de snelheden van beide voertuigen op $t_2$. Wie is dan het snelst?
        \end{hint}

        \begin{solution}\nl
            \underline{Gegevens:}\\
            $v_M(t) = 20 \, \mathrm{m/s} - 2{,}0 \, \mathrm{m/s^2} \cdot t$ en $v_V(t) = 8{,}0 \, \mathrm{m/s}$; $t_2 = 9{,}0 \, \mathrm{s}$\\

            \underline{Gevraagd:} welk voertuig het andere inhaalt\\

            \underline{Oplossing:}\\
            \[
                v_M(9{,}0 \, \mathrm{s}) = 20 \, \mathrm{m/s} - 2{,}0 \, \mathrm{m/s^2} \cdot 9{,}0 \, \mathrm{s} = 2{,}0 \, \mathrm{m/s}
            \]
            Dit is kleiner dan $v_V = 8{,}0 \, \mathrm{m/s}$. De vrachtwagen is dus sneller dan de motor en haalt hem in. Beide voertuigen rijden in dezelfde zin, dus ze kruisen elkaar niet in tegengestelde zin.
        \end{solution}
    \end{question}

    \begin{question}
        Wat is de grootste afstand tussen de motor en de vrachtwagen tijdens het tijdsinterval tussen $t_1$ en $t_2$?

        \[ \answer[tolerance=0.01,onlineshowanswerbutton]{9.0} \, \mathrm{m} \]

        \begin{hint}
            De afstand tussen beide voertuigen is eerst groeiend en daarna weer afnemend. Wat weet je over hun snelheden op het moment dat de afstand het grootst is?
        \end{hint}

        \begin{solution}\nl
            \underline{Gegevens:}\\
            $v_M(t) = 20 \, \mathrm{m/s} - 2{,}0 \, \mathrm{m/s^2} \cdot t$ en $v_V(t) = 8{,}0 \, \mathrm{m/s}$\\

            \underline{Gevraagd:} maximale afstand $x_M - x_V$\\

            \underline{Oplossing:}\\
            De afstand is maximaal als hij niet meer groeit en nog niet afneemt: dan zijn beide snelheden gelijk.
            \[
                20 \, \mathrm{m/s} - 2{,}0 \, \mathrm{m/s^2} \cdot t = 8{,}0 \, \mathrm{m/s} \quad \Rightarrow \quad t = 6{,}0 \, \mathrm{s}
            \]
            De posities op dat moment zijn:
            \begin{align*}
                x_M(6{,}0 \, \mathrm{s}) &= 20 \, \mathrm{m/s} \cdot 6{,}0 \, \mathrm{s} - 1{,}0 \, \mathrm{m/s^2} \cdot (6{,}0 \, \mathrm{s})^{2} = 84 \, \mathrm{m}\\
                x_V(6{,}0 \, \mathrm{s}) &= 27 \, \mathrm{m} + 8{,}0 \, \mathrm{m/s} \cdot 6{,}0 \, \mathrm{s} = 75 \, \mathrm{m}
            \end{align*}
            De grootste afstand is $84 \, \mathrm{m} - 75 \, \mathrm{m} = 9{,}0 \, \mathrm{m}$.
        \end{solution}
    \end{question}
\end{exercise}
%———————————————————————————————————————

%———————————————————————————————————————
% OEFENING 7 (moeilijkst): kromlijnige beweging, tangentiele en normale versnelling
\begin{exercise}
    Een bal wordt weggeschopt. Zijn positie wordt beschreven ten opzichte van een horizontale $x$-as (in de bewegingsrichting) en een verticale $y$-as (omhoog) door
    \begin{align*}
        x(t) &= 5{,}0 \, \mathrm{\frac{m}{s}} \cdot t\\
        y(t) &= 2{,}0 \, \mathrm{m} + 12 \, \mathrm{\frac{m}{s}} \cdot t - 5{,}0 \, \mathrm{\frac{m}{s^2}} \cdot t^{2}
    \end{align*}
    met $t$ in seconden. Deze vergelijkingen gelden tot de bal de grond ($y = 0{,}0 \, \mathrm{m}$) raakt.

    \begin{question}
        Bereken de grootte van de snelheid van de bal op $t = 0{,}50 \, \mathrm{s}$.

        \[ v = \answer[tolerance=0.01,onlineshowanswerbutton]{8.6} \, \mathrm{m/s} \]

        \begin{hint}
            De snelheidsvector heeft een $x$- en een $y$-component. Hoe vind je elke component vanuit de positievergelijking?
        \end{hint}

        \begin{solution}\nl
            \underline{Gegevens:}\\
            $x(t) = 5{,}0 \, \mathrm{m/s} \cdot t$ en $y(t) = 2{,}0 \, \mathrm{m} + 12 \, \mathrm{m/s} \cdot t - 5{,}0 \, \mathrm{m/s^2} \cdot t^{2}$; $t = 0{,}50 \, \mathrm{s}$\\

            \underline{Gevraagd:} $v$\\

            \underline{Oplossing:}\\
            De componenten van de snelheid zijn de afgeleiden van de posities:
            \begin{align*}
                v_x(t) &= x'(t) = 5{,}0 \, \mathrm{m/s}\\
                v_y(t) &= y'(t) = 12 \, \mathrm{m/s} - 10 \, \mathrm{m/s^2} \cdot t
            \end{align*}
            Op $t = 0{,}50 \, \mathrm{s}$ is $v_y = 12 \, \mathrm{m/s} - 10 \, \mathrm{m/s^2} \cdot 0{,}50 \, \mathrm{s} = 7{,}0 \, \mathrm{m/s}$. De grootte van de snelheid is:
            \[
                v = \sqrt{v_x^{2} + v_y^{2}} = \sqrt{(5{,}0 \, \mathrm{m/s})^{2} + (7{,}0 \, \mathrm{m/s})^{2}} = 8{,}6 \, \mathrm{m/s}
            \]
        \end{solution}
    \end{question}

    \begin{question}
        Ontbind de versnelling van de bal op $t = 0{,}50 \, \mathrm{s}$ in een tangentiële component $a_t$ (evenwijdig met de snelheid) en een normale component $a_n$ (loodrecht op de snelheid). Een negatieve $a_t$ betekent dat de component tegengesteld is aan de snelheid.

        \[ a_t = \answer[tolerance=0.02,onlineshowanswerbutton]{-8.1} \, \mathrm{m/s^2} \qquad \text{en} \qquad a_n = \answer[tolerance=0.02,onlineshowanswerbutton]{5.8} \, \mathrm{m/s^2} \]

        \begin{hint}
            Bepaal eerst de versnellingsvector zelf, en teken daarna de snelheids- en versnellingsvector samen in één schets.
        \end{hint}
        \begin{hint}
            Welke hoek maken de snelheidsvector en de versnellingsvector met elkaar? De component van $\vec{\mathbf{a}}$ langs en loodrecht op $\vec{\mathbf{v}}$ volgt uit die hoek.
        \end{hint}

        \begin{solution}\nl
            \underline{Gegevens:}\\
            $v_x = 5{,}0 \, \mathrm{m/s}$ en $v_y = 7{,}0 \, \mathrm{m/s}$ op $t = 0{,}50 \, \mathrm{s}$\\

            \underline{Gevraagd:} $a_t$ en $a_n$\\

            \underline{Oplossing:}\\
            De componenten van de versnelling zijn de afgeleiden van de snelheidscomponenten:
            \begin{align*}
                a_x(t) &= v_x'(t) = 0 \, \mathrm{m/s^2}\\
                a_y(t) &= v_y'(t) = -10 \, \mathrm{m/s^2}
            \end{align*}
            De versnelling wijst dus recht naar beneden, met grootte $a = 10 \, \mathrm{m/s^2}$. De snelheidsvector maakt een hoek met de $x$-as van:
            \[
                \theta = \tan^{-1}\left(\frac{7{,}0 \, \mathrm{m/s}}{5{,}0 \, \mathrm{m/s}}\right) = 54{,}5^\circ
            \]
            \begin{center}
                \begin{tikzpicture}[scale=1]
                    \coordinate (P) at (0,0);
                    \coordinate (O) at (-3.2,-2.6);
                    \tekenassen[1.2]{O}{0}
                    \draw[snelheid] (P) -- (54.46:2.8) node[above right] {$\vec{\mathbf{v}}$};
                    \draw[versnelling] (P) -- (0,-2.4) node[below] {$\vec{\mathbf{a}}$};
                    \draw[versnellingcomponent] (P) -- (234.46:1.95) node[below left] {$\vec{\mathbf{a}}_{t}$};
                    \draw[versnellingcomponent] (P) -- (-35.54:1.4) node[right] {$\vec{\mathbf{a}}_{n}$};
                    \draw[gray, dashed, thin] (234.46:1.95) -- (0,-2.4) -- (-35.54:1.4);
                    \fill (P) circle (1.5pt);
                \end{tikzpicture}

                Schets (niet op schaal)
            \end{center}
            De versnelling wijst naar beneden ($-90^\circ$), dus de hoek tussen $\vec{\mathbf{v}}$ en $\vec{\mathbf{a}}$ is $\varphi = 90^\circ + 54{,}5^\circ = 144{,}5^\circ$. Daaruit volgt:
            \begin{align*}
                a_t &= a \cdot \cos(\varphi) = 10 \, \mathrm{m/s^2} \cdot \cos(144{,}5^\circ) = -8{,}1 \, \mathrm{m/s^2}\\
                a_n &= a \cdot \sin(\varphi) = 10 \, \mathrm{m/s^2} \cdot \sin(144{,}5^\circ) = 5{,}8 \, \mathrm{m/s^2}
            \end{align*}
            Controle: $\sqrt{(8{,}1 \, \mathrm{m/s^2})^{2} + (5{,}8 \, \mathrm{m/s^2})^{2}} \simeq 10 \, \mathrm{m/s^2}$. De tangentiële component is negatief: de bal vertraagt op dat moment.
        \end{solution}
    \end{question}

    \begin{question}
        Op welk tijdstip $t$ is de snelheid van de bal het kleinst, en hoe groot is die snelheid dan?

        \[ t = \answer[tolerance=0.01,onlineshowanswerbutton]{1.2} \, \mathrm{s} \qquad \text{en} \qquad v_{\min} = \answer[tolerance=0.01,onlineshowanswerbutton]{5.0} \, \mathrm{m/s} \]

        \begin{hint}
            De grootte van de snelheid hangt af van $v_x$ en $v_y$. Welke van beide componenten verandert er in de tijd, en wat is dus de kleinst mogelijke waarde ervan?
        \end{hint}

        \begin{solution}\nl
            \underline{Gegevens:}\\
            $v_x = 5{,}0 \, \mathrm{m/s}$ (constant) en $v_y(t) = 12 \, \mathrm{m/s} - 10 \, \mathrm{m/s^2} \cdot t$\\

            \underline{Gevraagd:} $t$ en $v_{\min}$\\

            \underline{Oplossing:}\\
            De snelheid is $v = \sqrt{v_x^{2} + v_y^{2}}$. Omdat $v_x$ constant is, is $v$ minimaal als $\lvert v_y \rvert$ minimaal is, dus als $v_y = 0$:
            \[
                12 \, \mathrm{m/s} - 10 \, \mathrm{m/s^2} \cdot t = 0 \quad \Rightarrow \quad t = 1{,}2 \, \mathrm{s}
            \]
            De snelheid is dan enkel nog horizontaal: $v_{\min} = v_x = 5{,}0 \, \mathrm{m/s}$.

            Op dat moment staat $\vec{\mathbf{a}}$ (verticaal) loodrecht op $\vec{\mathbf{v}}$ (horizontaal). Er geldt dus $a_t = 0$ en $a_n = 10 \, \mathrm{m/s^2}$: de snelheid verandert enkel van richting, niet van grootte.
        \end{solution}
    \end{question}

    \begin{question}
        Op welke horizontale positie $x$ raakt de bal de grond?

        \[ x = \answer[tolerance=0.01,onlineshowanswerbutton]{12.8} \, \mathrm{m} \]

        \begin{hint}
            Welke waarde heeft $y$ op het moment dat de bal de grond raakt? Los eerst op naar de tijd.
        \end{hint}

        \begin{solution}\nl
            \underline{Gegevens:}\\
            $x(t) = 5{,}0 \, \mathrm{m/s} \cdot t$ en $y(t) = 2{,}0 \, \mathrm{m} + 12 \, \mathrm{m/s} \cdot t - 5{,}0 \, \mathrm{m/s^2} \cdot t^{2}$\\

            \underline{Gevraagd:} $x$ als $y = 0$\\

            \underline{Oplossing:}\\
            We stellen $y(t) = 0$ (eenheden even weggelaten, $t$ in $\mathrm{s}$):
            \[
                5t^{2} - 12t - 2 = 0 \quad \Rightarrow \quad t = \frac{12 \pm \sqrt{144 + 40}}{10} = \frac{12 \pm 13{,}56}{10}
            \]
            Enkel de positieve oplossing heeft fysische betekenis: $t = 2{,}56 \, \mathrm{s}$. Dan is:
            \[
                x = 5{,}0 \, \mathrm{m/s} \cdot 2{,}56 \, \mathrm{s} = 12{,}8 \, \mathrm{m}
            \]
        \end{solution}
    \end{question}
\end{exercise}
%———————————————————————————————————————

\end{document}